\section{Finite approximations and cap selection}\label{sec:cap}

Retain the first \(N\) annuli of the datum in \eqref{eq:datum}, and use a
common collision cutoff \(b_N\to\infty\) in \eqref{eq:colored}. All assertions
in this section first concern fixed \(N,b_N\). Constants in the cold
estimates and in the caps will be uniform in \(N\).

Choose once and for all a nonnegative smooth function
\(\chi\in C_c^\infty(\R^3_x\times\R^3_v)\), not identically zero.
After reducing \(s_0\), its periodic copies at spatial scale \(s_N\) are
disjoint initially. The hot initial density is
\[
h(0,x,v)=\lambda s_N\chi(x/s_N,v),
\]
periodized in \(x\). The cold initial densities are independent of
\(\lambda\). In particular the norm
\(\|\chi(x/s_N,v)\|_H\) is a fixed positive number independent of \(N\).

\begin{proposition}\label{prop:finite-cap}
There is a fixed \(T_0>0\) such that the following statements hold.
\begin{enumerate}[label=\textup{(\alph*)}]
\item With \(\lambda=0\), the finite approximations exist on \([0,T_0]\)
and satisfy \(0\le h\le t^2H\).
\item There is a fixed finite \(L_0\) such that, for every
\(0<T\le T_0\) and every \(N\), some \(\lambda\in[0,L_0]\) gives an
approximation on \([0,T]\) with
\begin{equation}\label{eq:cap}
0\le h(t,x,v)\le(t+s_N)H(v),\qquad
\max_{0\le t\le T}\frac{\|h(t)\|_H}{t+s_N}=1.
\end{equation}
\end{enumerate}
In both cases all cold components obey Proposition~\ref{prop:cold-control}.
\end{proposition}

\begin{proof}
We first give the fixed-approximation local theory used in the selection.
Lemma~\ref{lem:hot-bound} bounds the gain bilinearly in \(\|\cdot\|_H\).
For the cutoff loss,
\[
\|\nu_f^{b_N}\|_\infty
 \le2\pi b_N\esssup_x\int_{\R^3}|f(x,w)|\dd w
 \le2\pi b_N\|f\|_H\int_{\R^3}H(w)\dd w,
\]
so gain and loss are locally Lipschitz on bounded
balls of the weighted norm. Free transport is an isometry in that norm.

Picard iteration for the integral equation along characteristics gives
a unique local solution in the space of bounded measurable functions
with finite weighted norm. An integrating factor for the loss, followed
by positive iteration of the gain, gives nonnegativity. One can equivalently
perform this iteration for the labeled system, using the norm of the sum
of the absolute values of all components, and for differences the sum of
the absolute differences. The quadratic estimates close in that norm.
The Gaussian spatial tails in Proposition~\ref{prop:cold-control} make
the sums over periodic copies finite in the relevant norm. Periodic
covariance is preserved by uniqueness.

The solution continues as long as its total weighted norm remains
bounded. It depends continuously on its initial datum, uniformly on
a fixed interval on which that norm is bounded. These assertions follow
from the same integral contraction and its difference estimate on short
subintervals. They do not require norm continuity of spatial translations:
one works with the integral equation along the free paths.
In particular,
\[
\|h(t)-h(s,x-(t-s)v,v)\|_H\le C_{N,b_N}|t-s|
\]
on a bounded solution interval. Since transport is an isometry,
\(t\mapsto\|h(t)\|_H\), and hence \(t\mapsto\|h(t)\|_H/(t+s_N)\),
are continuous.

Let \(z(t)=\|h(t)\|_H\). Dropping the loss in the hot equation,
the cold and kernel estimates give its scalar integral comparison with
\begin{equation}\label{cap:scalar}
z'(t)\le
 \left(\frac{C_1}{t+s_N}+C_1\right)z(t)+C_1z(t)^2+C_A t^{A-1}.
\end{equation}
Here \(A\) can be any sufficiently large fixed number, and its source
constant is uniform in \(N\). This is an upper integral inequality;
no differentiability of the essential supremum is needed.

Fix \(A>C_1+2\). Choose \(D>0\) so that
\((A-C_1)D>C_A\). Since \(t/(t+s_N)\le1\), the function \(Dt^A\)
is a supersolution of \eqref{cap:scalar} after shortening a fixed
preliminary time interval to absorb the bounded linear and quadratic
terms. Shorten it again to ensure \(Dt^A\le t^2\). Comparison from
zero proves (a). The cold bound is independent of \(h\), so it also
ensures continuation for each fixed \(N,b_N\) up to this time.

For (b), shorten \(T_0\) to be less than one, so that the zero-seed
solution satisfies a strict cap:
\[
\frac{\|h(t)\|_H}{t+s_N}\le\frac{t^2}{t+s_N}\le T_0<1.
\]
Choose \(L_0\) so large that
\(L_0\|\chi(x/s_N,v)\|_H>1\). The cap then fails at time zero for
\(\lambda=L_0\).

For fixed \(N,b_N,T\), let \(\mathcal O\) be the set of seed parameters
in \([0,L_0]\) whose solution exists up to \(T\) and obeys a strict cap
there. Continuous dependence and continuation make \(\mathcal O\)
relatively open. Its component containing zero has a right endpoint
\(\lambda_*\) strictly below \(L_0\): indeed a neighborhood of \(L_0\)
already fails the initial cap. Approach \(\lambda_*\) from within that
component. The cap and the fixed cold bound give a common finite
weighted norm on \([0,T]\), with constants permitted to depend on
\(N,b_N\). The local difference estimates on finitely many short
subintervals give a limit solution at \(\lambda_*\), with cap at most
one. It continues to \(T\). If its maximum cap ratio were less than
one, continuous dependence would extend the component past
\(\lambda_*\), a contradiction. This proves \eqref{eq:cap}.
\end{proof}

\begin{remark}\label{cap:parameter-order}
The coefficient \(C_1\) may be large, in particular when
\(\alpha\gamma\) is small. It is fixed before \(A,D,T_0\) are chosen.
The arbitrary-power tail estimates are available on a preliminary
short interval; their constants can be fixed there before the final
shortening. Thus the choice does not depend on \(N\), and imposes
no restriction on the later choice of shrinking horizons.
\end{remark}
